\paper
{Development of a Deep Learning-Based NIR Spectroscopic System for Plastic Classification}
{Vishishta Dilsara, Thisal Vitharana, Rukshika Fernando, Vindya Senanayake, Dhammika Weerathunga}
{Vishishta Dilsara$^{1}$, Thisal Vitharana$^{2}$, Rukshika Fernando$^{2}$, Vindya Senanayake$^{1}$, Dhammika Weerathunga$^{2}$}
{$^{1}$Faculty of Computing, University of Sri Jayewardenepura, Sri Lanka\\
	$^{2}$Faculty of Applied Sciences, University of Sri Jayewardenepura, Sri Lanka\\
	Email: vishishtadilsara2002@gmail.com, as2022159@sci.sjp.ac.lk, as2022376@sci.sjp.ac.lk, vindya@sjp.ac.lk, dhammikaweerathunga@sjp.ac.lk}
{The rising levels of plastic waste have heightened the demand for automated and reliable classification of plastics to facilitate effective recycling. Near-Infrared (NIR) spectroscopy technology has gained prominence as a widely used technique owing to its fast, non-destructive nature and material-specific spectral response. However, classifying plastics based on their NIR spectra poses great difficulty owing to factors such as spectral noise, variability, and overlapping absorption regions caused by similar molecular bonds, e.g., C-H, O-H, N-H, especially when comparing plastics having chemically similar structures such as PET/PVC and HDPE/LDPE. The conventional machine learning and chemometric approach to classify plastics requires significant amounts of manual data pre-processing (such as normalization and smoothing). A comprehensive overview of deep learning techniques, especially the use of CNNs, is performed in this study to determine how effectively these techniques can be used to classify plastics by using NIR spectral information. State-of-the-art deep learning methods for this specific task have been assessed based on several performance evaluation metrics, such as accuracy, precision, recall, F1 score, and other measures. There are some drawbacks in this area, including poor interpretability of models, poor generalization ability to different datasets, and computational complexity for implementing models on mobile devices. The results of this study identify these limitations and suggest directions for further research.
	
	\textbf{Keywords:} Near-Infrared Spectroscopy, Convolutional Neural Networks, Plastic Waste Identification, Spectral Classification, Deep Learning
}